\subsection{Systemprüfung und Installation}
\label{subsec:systempruefung-installation}

Der read-only ausgelegte \texttt{doctor} prüft den laufenden
Python-Interpreter, Node.js, npm und npx sowie das optionale \texttt{uv}, lädt
Profil und wirksame Konfiguration und kontrolliert relevante Projektpfade,
Abhängigkeiten und Ports. Deaktivierte Schichten gelten nicht als Fehler.
Cloudprofile ergänzen Containerprüfungen, wobei fehlendes Docker im allgemeinen
Doctor nur warnt; erst \texttt{container doctor} behandelt es als
Voraussetzung. Datenbankverbindung und native Desktop-Bibliotheken besitzen mit
\texttt{db doctor --connect} beziehungsweise \texttt{tauri doctor} bewusst
separate, vertiefende Diagnosewege. Letzterer prüft zusätzlich Rust/Cargo,
Tauri-CLI und plattformspezifische Bibliotheken
\parencite{kleiveist2026tooling,kleiveist2026configuration}.

\texttt{install} verwendet für das Frontend bei vorhandenem Lockfile
\texttt{npm ci}, andernfalls \texttt{npm install}. Für ein aktives Backend
legt es \path{backend/.venv} an und installiert die Basisanforderungen sowie
profilabhängig Datenbank- und PostgreSQL-Ergänzungen; \texttt{uv} ist ein
optionaler schneller Pfad, für den ein pip/venv-Fallback existiert. Ferner wird
bei Bedarf eine kleine Tooling-Umgebung für Pytest eingerichtet. Chromium wird
nur installiert, wenn Playwright im Projekt konfiguriert ist. Deaktivierte
Frontend- oder Backendschichten werden übersprungen, und eine lokale
\texttt{.env} wird weder erzeugt noch verändert. Native Betriebssystem- und
Rust-Voraussetzungen bleiben Aufgabe von \texttt{tauri install}. Damit bündelt
der Standardweg die relevanten npm- und Python-Schritte, ohne jeden
Sonderworkflow vorwegzunehmen.

Die README nennt Git, Python ab Version~3.11 und Node.js ab Version~20 als
Grundvoraussetzungen sowie Rust Stable für Desktopprofile. Der allgemeine
Doctor prüft jedoch weder Git noch diese Versionsbereiche; auch der Tauri-Doctor
meldet jeden vorhandenen aktiven Rust-Toolchain als gültig. Die festgelegten
Python- und Node-Hauptversionen werden dagegen in den CI-Workflows explizit
eingerichtet. Dies ist eine Grenze zwischen dokumentierter Voraussetzung und
lokaler Durchsetzung, nicht eine zusätzliche Tooling-Funktion
\parencite{kleiveist2026template,kleiveist2026ci}.
